\documentclass[border=4pt]{standalone}
\usepackage[siunitx]{circuitikz}

\begin{document}
\begin{circuitikz}[european resistors]
  % Illustrative DC input-error equivalent for an inverting stage. The
  % offset source and bias currents are signed model quantities, not a
  % transistor-level or device-specific internal schematic.
  \node[font=\small\bfseries] at (6.0,9.45)
    {Illustrative inverting-stage DC input-error equivalent};
  \node[font=\small] at (6.0,8.90)
    {(sign-declared teaching model---not device-specific)};

  \node[op amp, scale=1.15, anchor=center] (U1) at (8.8,4.55) {};
  \coordinate (JM) at (6.15,0 |- U1.-);
  \coordinate (OUT) at (10.55,0 |- U1.out);
  \coordinate (STH) at (0.70,0 |- U1.-);

  % Signal-source Thevenin equivalent and the explicit inverting input
  % resistor. I_B- is positive into the physical op-amp input pin.
  \draw (0.70,0) node[ground] {}
    to[V, l={$V_{\mathrm{TH}}$}] (STH)
    to[R, l={$R_{\mathrm{TH}}$}] (3.05,0 |- STH)
    to[R, l={$R_{\mathrm{IN}}$}] (JM);
  \draw (JM) to[short, i>^={$I_{B-}$}] (U1.-);
  \draw (JM) -- (6.15,7.65)
    to[R, l={$R_F$}] (10.55,7.65)
    -- (OUT);
  \node[circ] at (JM) {};
  \node[font=\small, anchor=south east] at (JM) {$v_-$};

  % Non-inverting reference Thevenin model, optional bias-current
  % compensation resistor, and the signed series input-offset source.
  \coordinate (RTHP) at (2.10,0 |- U1.+);
  \coordinate (RREFEND) at (3.70,0 |- U1.+);
  \coordinate (RBEND) at (5.10,0 |- U1.+);
  \coordinate (OSOUT) at (6.35,0 |- U1.+);
  \draw (2.10,0) node[ground] {}
    to[V, l_={$V_{\mathrm{REF,TH}}$}] (RTHP)
    to[R, l={$R_{\mathrm{REF,TH}}$}] (RREFEND)
    to[R, l={$R_B$}] (RBEND)
    to[V, l_={$V_{\mathrm{OS}}$}] (OSOUT);
  \draw (OSOUT) to[short, i>_={$I_{B+}$}] (U1.+);
  \node[font=\small] at (5.25,4.58) {$-$};
  \node[font=\small] at (6.20,4.58) {$+$};
  \node[font=\small, anchor=south east] at (U1.+) {$v_+$};

  % Output, finite load, supplies, and the declared circuit reference.
  \draw (U1.out) -- (OUT)
    -- (11.75,0 |- OUT)
    node[ocirc, label={[font=\small]right:{TPO: $v_o$}}] {};
  \draw (OUT) to[R, l_={$R_L$}] (10.55,0)
    node[ground, label={[font=\small]right:{TP0: $0$ V}}] {};
  \draw (U1.up) -- (U1.up |- 0,6.15) node[vcc] {$V_{S+}$};
  \draw (U1.down) -- (U1.down |- 0,2.85) node[vee] {$V_{S-}$};
  \node[circ] at (OUT) {};
  \node[font=\small, anchor=west] at (9.35,5.65) {$U_1$};

  % Explicit sign and first-order compensation declarations.
  \node[draw, fill=white, align=center, font=\small] at (6.0,-1.15)
    {$V(\mathrm{OS+})-V(\mathrm{OS-})\equiv V_{\mathrm{OS}}$;
     \quad $I_{B+},I_{B-}>0$ into the input pins};
  \node[draw, fill=white, align=center, font=\small] at (6.0,-2.05)
    {with independent voltage sources set to zero:
     \quad $R_{\mathrm{REF,TH}}+R_B
       \approx R_F\parallel(R_{\mathrm{IN}}+R_{\mathrm{TH}})$};
\end{circuitikz}
\end{document}
